\documentclass[../DoAn.tex]{subfiles}
\begin{document}

\begin{center}
    \Large{\textbf{ABSTRACT}}\\
\end{center}
\vspace{1cm}
The demand for learning Japanese in Vietnam has grown steadily for many years, yet learning this language requires frequent exposure to authentic audio-visual content and active recall of newly acquired knowledge. Popular platforms such as Duolingo, WaniKani, and YouTube video channels do not fully address this problem: they either lack structured video content, lack a knowledge-check mechanism immediately following each lesson, or do not support the diverse question types required for Japanese, and are not localized for Vietnamese learners.

To address this problem, this thesis builds an online Japanese learning platform that combines lecture videos with an interactive quiz immediately following each video, using a web architecture that separates the frontend and backend through a RESTful API — an approach that allows the two components to be developed independently while remaining easy to extend and maintain for a project built by a single developer.

The system's backend is built with Spring Boot and a MySQL database that manages all business logic for courses, lessons, quizzes, learning progress, and JWT-based user authentication; the frontend uses React together with TanStack Query and Zustand for data fetching and UI state management. Content is organized into a Level – Course – Lesson – Quiz hierarchy, allowing students to enroll in courses (free or paid via VNPay), watch lecture videos with progress tracking, and complete a 10-question quiz divided into three question-type groups (vocabulary, content, and sequence) immediately after finishing the video. Learning results, attempt history, and daily learning streaks are stored so students can track their progress; administrators have a separate interface to manage courses, lessons, and the question bank.

The thesis's main contribution is a complete, deployed, and tested web system in which the quiz subsystem is generalized around a question-type model, allowing new question types to be added in the future without changing the database schema, together with a clear separation of access control between students and administrators and a JWT authentication scheme with refresh-token rotation for user security.

\end{document}